% ---------------------------------------------------------------------------
% §3 EMPIRICAL STRATEGY — paste-ready revize.
% Yapı (Fuat feedback s.8): Empirical Strategy üst başlık; girişte iki aşama
% + yol haritası; sonra Data, Calibration Outcomes, Persona Construction,
% Calibration Procedure, Experimental Design.
% TODO'lar: (1) \citep{zhang2025verbalized} için bib girdisi ekle
% (Verbalized Sampling makalesi). (2) \ref{sec:results} etiketini kendi
% Results section etiketinle değiştir. (3) Distributional-alignment gerekçe
% cümlelerine Fuat'ın referanslı metni gelecek (kendisinden alınacak).
% ---------------------------------------------------------------------------

\section{Empirical Strategy}
\label{sec:empirical-strategy}

This study proceeds in two stages. The first stage calibrates and selects
the inference protocol used to generate survey responses from persona-based
synthetic agents, which we refer to as digital twins. We use two held-out
TGSS outcomes as ground truth and vary persona content, sampling strategy,
temperature, address mode, model, and prompt framing. Each setup is
evaluated at both the distributional and respondent levels. The aim is to
identify the protocol that most closely reproduces observed survey
responses before applying it in the experimental stage.

The second stage applies the calibrated inference protocols in a survey
experiment conducted on the synthetic population. Each digital twin is
exposed, in independent prompt turns, to two versions of the same protest
event that differ only in the surrounding political structure: a threat
condition, set against an active cross-border military operation, and an
opportunity condition, set during an ongoing negotiation process. This
within-persona design allows us to compare responses from the same
synthetic respondent across the two conditions. Treatment effects are
evaluated as aggregate distributional shifts, consistent with the primary
calibration target established in the first stage.

The remainder of this section introduces the data, the calibration outcomes
and their evaluation metrics, the construction of the personas, the
calibration procedure, and the experimental design, in that order.

\subsection{Data}
\label{sec:data}

This study draws on the 2024 Turkish General Social Survey (TGSS 2024), a
nationally representative face-to-face survey of 2,615 adults conducted
across Turkey's twelve NUTS-1 regions \citep{nisanci_tgss_2025}. The TGSS
covers demographic characteristics, socioeconomic position, political
attitudes and behavior, social trust and networks, religiosity, well-being,
media consumption, and civic participation. This breadth of
individual-level information makes the survey especially suitable for
persona-based simulation, since it allows us to construct respondent
profiles that go beyond basic demographics and incorporate social,
political, and attitudinal characteristics.

\subsection{Calibration Outcomes and Evaluation}
\label{sec:outcomes}

The calibration outcomes are selected to differ in response structure while
remaining substantively connected to the downstream experiment.
\textit{Pacdemons} (\textit{n} = 1,707) is a rare binary measure of
participation in street demonstrations during the previous twelve months,
with affirmative responses accounting for approximately five percent of
valid cases.\footnote{The TGSS uses a rotating question design, so
\textit{pacdemons} was administered to a subset of respondents (\textit{n} =
1,707) rather than to the full sample.} \textit{Womenwork} is a five-point
Likert item measuring agreement with the statement that family life suffers
when a woman works full time (\textit{n} = 2,588). Together, these outcomes
test whether the optimal inference setup remains stable across response
formats or requires adaptation to the measurement structure of the target
outcome.

The two outcomes also capture dimensions relevant to the downstream
experiment on protest for Kurdish-language education. Protest participation
measures the respondent's prior relationship to collective action, while
attitudes toward women's full-time work capture responses to a contested
question of social equality and social change. Both outcomes are excluded
from persona construction at every stage. Agents therefore never receive
the answers they are later asked to reproduce, and evaluation compares each
simulated response with the respondent's withheld survey response.

The primary calibration target is distributional alignment because the
downstream experiment compares aggregate response distributions across
treatment conditions. We therefore use metrics tailored to the measurement
scale of each outcome. For the binary protest-participation outcome, we
report Jensen--Shannon divergence (JSD) and the raw prevalence difference,
$\Delta p$. JSD measures overall dissimilarity between the observed and
simulated response distributions, while $\Delta p$ captures the direction
and magnitude of bias in positive-class prevalence
\citep{lin1991,bisbee2024}. For the five-point Likert outcome, we use the
Wasserstein-1 distance because it respects the ordered structure of the
response categories and penalizes discrepancies in proportion to the
distance between categories \citep{rubner2000}. We additionally report JSD
and comparisons of item-level means and standard deviations as descriptive
diagnostics, since recent research on synthetic survey respondents
documents underdispersion and variance compression in LLM-generated
responses \citep{bisbee2024,toubia2025}.

Because the synthetic population is generated at the respondent level
before responses are aggregated, we evaluate calibration performance at
both the aggregate and respondent levels. For \textit{pacdemons},
respondent-level performance is assessed using the Matthews correlation
coefficient, macro-F1, protest-positive recall (recall on the rare
affirmative class), and Cohen's kappa
\citep{cohen1960,matthews1975,powers2011}. For \textit{womenwork}, we
report Cohen's kappa and quadratic weighted kappa. Quadratic weighted kappa
is appropriate for this outcome because it distinguishes adjacent-category
errors from larger ordinal discrepancies on the five-point scale
\citep{cohen1968}. Final protocol selection considers distributional fit
and respondent-level agreement jointly. The selected protocol is therefore
the one that provides the strongest overall performance across both
dimensions.

\subsection{Persona Construction}
\label{sec:persona-construction}

The first decision in building a digital twin is how much of a respondent
to encode. A persona may include only basic demographics or a richer
profile spanning beliefs, economic position, social values, and political
attitudes. Because more information is not necessarily better, overlong
profiles may dilute task-relevant information or encourage reliance on
broad attribute associations. We therefore treat persona detail as one
dimension of the calibration design.

We begin from a curated set of roughly 117 TGSS variables, organized into
six thematic groups, namely demographics, belief and ideology, economic
position and attitudes, social values and attitudes, social-psychological
dispositions, and political values and attitudes.\footnote{The full list of
variables is reported in Appendix~\ref{app:variables}.}

Each variable is converted into a short natural-language line that pairs a
compact label with the respondent's value. Labels are generated once by
shortening each survey question into a two-to-four word form, then fixed
and reused for all respondents. Variables appear under group headings in a
fixed order, from demographics to political values and attitudes. The
result is a structured respondent profile built deterministically from the
survey record, with no generation step before inference.

We use a structured format as the baseline representation of each persona
rather than automatically converting survey records into continuous prose.
This format preserves every included variable, maintains a consistent
ordering across respondents, and avoids introducing uncontrolled variation
during persona construction. It therefore provides a standardized starting
point for evaluating differences across inference settings. As part of the
calibration procedure, we also test whether rewriting the shortlisted
personas into natural-language profiles improves performance relative to
this structured baseline.

The design produces eleven persona-content configurations, labeled C0
through C10, which specify how much of the respondent's survey record
enters the prompt. C0 is the demographic-only baseline, C1 is the full
profile, and C2 through C10 add or remove different combinations of
ideological, social, economic, psychological, and political information. By
varying these blocks, we test which forms of respondent information improve
calibration performance.\footnote{The exact composition of each
configuration is reported in Appendix~\ref{app:configs}, and
Appendix~\ref{app:persona-example} provides an anonymized example of a
persona profile used in the calibrated \textit{womenwork} protocol,
together with the full prompt in which it is embedded.} Applied to the full
TGSS sample, the design yields $2{,}615 \times 11 = 28{,}765$ distinct
persona profiles.

\subsection{Calibration Procedure}
\label{sec:calibration-procedure}

Having fixed how personas are built, we calibrate the inference setup that
converts each persona into a survey answer. We conduct this search
separately for the two calibration outcomes. The full grid of persona
configuration, sampling strategy, temperature, address mode, model, and
prompt presentation spans roughly 2,600 cells per outcome, and evaluating
every cell on the full sample would require more than eleven million model
completions before any multi-seed replication --- well beyond a defensible
computational budget. We therefore use a sequential funnel
(Figure~\ref{fig:calibration-funnel}), in which each step narrows the
candidate set by varying one component of the setup while holding the
remaining components fixed. The order follows the structure of the problem:
we begin with persona content, then vary inference behavior, model family,
and prompt presentation.\footnote{Full cell-level results and
implementation details are reported in
Appendix~\ref{app:calibration-results}.}

We first hold the inference setup fixed at direct answering, temperature
zero, first-person address, and GPT-4o-mini, then screen the eleven persona
configurations across both outcomes. This step isolates the effect of
persona content and reduces the full set of configurations to a smaller
candidate set for subsequent tests.

On the shortlisted configurations, we compare direct answering with
verbalized sampling across three temperatures, $T \in \{0, 0.4, 0.8\}$.
Direct answering records the model's selected option. Verbalized sampling
asks the model to produce a probability distribution over the valid
response options, from which we draw one response using a fixed
respondent-level seed \citep{zhang2025verbalized}. % TODO: bib girdisi ekle
Sampling strategy and temperature are varied jointly because both govern
the same stochastic layer of the setup: temperature shapes the variability
of what the model generates, while the sampling strategy determines whether
that variability is expressed as a single selected option or as an explicit
probability distribution. This design separates probabilistic answer
generation from the final sampled response and allows us to test whether
deterministic or probabilistic answering better recovers each
outcome.\footnote{The parsing procedure and example prompt are reported in
Appendix~\ref{app:verbalized-prompt}; stochastic seed construction is
reported in Appendix~\ref{app:stability}.}

We then re-run the leading per-outcome setups on three additional models
from different providers: GPT-5.4-mini (OpenAI), Gemini 2.5 Flash Lite
(Google), and Llama 3.3 70B (Meta). This step assesses whether the main
calibration patterns are specific to one model implementation or persist
across alternative model families. The four backbones are chosen because
they are among the most widely used commercial and open-weight model
families in applied synthetic-survey work, spanning three providers; this
anchors the robustness check in the models practitioners actually deploy
rather than in an exhaustive sweep of the model space.

Finally, we test whether the selected protocols are robust to plausible
changes in how the persona and the survey prompt are presented, holding the
underlying respondent information fixed. Four variants are considered:
switching the grammatical address of the prompt from first to second
person, inserting an instruction to reason briefly before selecting an
answer, prepending a short background paragraph describing the structure of
political ideology in Turkey,\footnote{The full Turkish text of the
background paragraph is reproduced in Appendix~\ref{app:background}.} and
rewriting the structured persona profile as a natural-language biography
that preserves every recorded fact while adding none.\footnote{The
rewriting procedure and full results of this ablation are reported in
Appendix~\ref{app:natural}.} Each variant is a choice a practitioner could
defensibly make when designing the simulation. If simulated responses
shifted under these variants, calibration results would rest on arbitrary
presentation decisions rather than on persona content; the first three
tests confirm that they do not, while the fourth shows that the
\emph{format} of the persona is not arbitrary at all: aggregate
distributional fit survives the biographical rewrite, but respondent-level
agreement degrades sharply, which is why the structured format is retained.

At the end of the funnel, we select one calibrated inference setup for each
outcome. These setups are fixed before the downstream survey experiment, so
experimental results cannot influence the choice of persona configuration,
sampling strategy, model, temperature, or prompt presentation. After
protocol selection, we assess stochastic stability across random seeds as a
robustness check.\footnote{Details are reported in
Appendix~\ref{app:stability}.} The results of the calibration funnel and
the selected protocols are reported in the next section.

% Figür bölümün SONUNA taşındı (Fuat feedback s.10): okuyucu tüm kavramları
% tanıdıktan sonra funnel'ı görür. [t] float olduğu için sayfa üstüne
% yerleşecektir; kaynaktaki konumu son sayfalara denk getirmek yeterli.
\begin{figure*}[t]
\centering
\begin{tikzpicture}[
  boxstep/.style={
    rectangle, rounded corners=4pt,
    draw=black!60, fill=gray!10,
    text width=40mm, align=center,
    minimum height=28mm,
    inner ysep=6pt, inner xsep=5pt
  },
  boxfinal/.style={
    rectangle, rounded corners=4pt,
    draw=black!80, fill=black!15,
    text width=40mm, align=center,
    minimum height=28mm,
    inner ysep=6pt, inner xsep=5pt
  },
  flow/.style={-{Stealth[length=5pt,width=4pt]}, thick, black!50}
]
\node[boxstep] (s1) at (0,0) {%
  \textbf{\small 1.\ Configuration screening}\\[4pt]
  \scriptsize 11 configs $\times$ 2 outcomes;
  baseline: direct answering,
  $T{=}0$, first-person framing
};
\node[boxstep] (s2) at (50mm,0) {%
  \textbf{\small 2.\ Sampling \& temperature}\\[4pt]
  \scriptsize Shortlisted configs $\times$
  \{direct, verbalized\}
  $\times\; T \in \{0,\;0.4,\;0.8\}$
};
\node[boxstep] (s3) at (100mm,0) {%
  \textbf{\small 3.\ Cross-model robustness}\\[4pt]
  \scriptsize Leading setups re-run across:
  GPT-4o-mini, GPT-5.4 mini,
  Gemini 2.5 Flash Lite, Llama 3.3 70B
};
\node[boxstep] (s4) at (0,-48mm) {%
  \textbf{\small 4.\ Presentation stress tests}\\[4pt]
  \scriptsize Address mode, explicit
  reasoning instruction, ideology
  background paragraph, natural rewrite
};
\node[boxfinal] (sf) at (50mm,-48mm) {%
  \textbf{\small Calibrated protocol}\\[4pt]
  \scriptsize One inference setup per outcome,
  fixed before the main experiment
};
%% Arrows
\draw[flow] (s1.east) -- (s2.west);
\draw[flow] (s2.east) -- (s3.west);
\draw[flow] (s3.south) -- ++(0,-10mm) -- ++(-100mm,0) -- (s4.north);
\draw[flow] (s4.east) -- (sf.west);
\end{tikzpicture}
\caption{Sequential calibration funnel. Each step narrows the candidate set
by evaluating a different component of the inference setup. Final protocol
selection is outcome-specific: for the rare binary outcome, we prioritize
respondent-level classification performance; for the ordered Likert
outcome, we jointly evaluate aggregate ordinal distributional fit and
respondent-level ordinal agreement. Stochastic stability across seeds is
assessed after protocol selection and reported in
Appendix~\ref{app:stability}.}
\label{fig:calibration-funnel}
\end{figure*}

\subsection{Experimental Design}
\label{sec:experimental_design}

The experiment asks how the political structure surrounding a protest event
shapes public reactions to it. We manipulate a single factor with two
levels, which we label \emph{political structure} following the
opportunity--threat framework of the political process tradition: a
\textbf{threat} condition, in which the protest occurs during an active
security episode, and an \textbf{opportunity} condition, in which the same
protest occurs during an ongoing peace and negotiation process. Each
condition is operationalized as a short news-style vignette. The threat
vignette reports a Ministry of National Defense statement on intensified
armed clashes in northern Syria, including casualties and heightened
security measures. The opportunity vignette reports ongoing talks between
the government and Kurdish political representatives, widely attended
Newroz celebrations, and Newroz messages issued by the President and party
leaders. Both vignettes close with an identical description of the focal
event: a march in Ankara at which a group voices the demand for
Kurdish-language education.\footnote{The full Turkish texts of both
vignettes, together with English translations, the original item wordings,
and the prompt templates, are reported in
Appendix~\ref{app:experimental-materials}.}

Two survey items follow each vignette and are given equal analytical
weight. \emph{Movement legitimacy} asks respondents to rate agreement with
the statement that the group is waging a legitimate rights struggle, on a
five-point Likert scale. \emph{Behavioral intent} asks whether the
respondent would personally attend such a march, as a binary choice. The
two items capture complementary faces of public response, an evaluative
judgment about the movement and a costly behavioral disposition, and mirror
the response structures on which the inference protocols were calibrated:
an ordered attitudinal item and a rare binary participation item.

Each outcome is administered under the protocol calibrated for its response
format (Section~\ref{sec:results}). % TODO: Results section etiketini kontrol et
The Likert legitimacy item uses the C6 persona configuration with
verbalized sampling at $T=0.8$; the binary intent item uses C6 with direct
answering at $T=0$. No information beyond the persona profile, the
vignette, and the survey item enters the prompt. The prompt templates are
identical to the calibration stage except for the inserted vignette.

The design is within-persona: all 2{,}615 twins respond under both
political structures, in independent prompt turns with no shared
conversational state, so that responses in one condition cannot contaminate
the other. Condition-specific respondent-level seeds control the stochastic
sampling step. The primary estimand is the within-persona paired contrast
between the threat and opportunity conditions, aggregated over the
synthetic population; we report paired $t$-tests with standardized effect
sizes ($d_z$) at the population level and within subgroups. Because each
twin serves as its own control, the design removes all between-respondent
heterogeneity from the contrast, which is precisely the counterfactual that
observational protest research cannot recover: the same public observed
under two political structures.

Consistent with the calibration results, the experiment is interpreted at
the level of aggregate and subgroup response distributions. The protocols
were validated for distributional resemblance, and the experimental claims
inherit that scope.
